\pagenumbering{arabic} % a Kivonattól kezdve folyamatos arab számozás
\setcounter{page}{1}

\selecthungarian

%----------------------------------------------------------------------------
% Abstract in Hungarian
%----------------------------------------------------------------------------
\chapter*{Kivonat}\addcontentsline{toc}{chapter}{Kivonat}

A különféle közösségek és diákkörök mindennapi működésük, belső rendezvényeik és táboraik során számos költségmegosztási és készletkezelési feladattal szembesülnek: a közösségi készletek felhasználását, a tagok közös kiadásait, az előzetes költség-hozzájárulásokat, valamint a csoportos kiadásokat nyomon kell követni. Ezt a belső elszámolást a csoportok jellemzően papíron, táblázatkezelőben vagy egyszerű feljegyzésekkel vezetik. Ezek a módszerek nehezen átláthatók, hibalehetőségeket hordoznak, nem biztosítanak valós idejű egyenlegkövetést a tagoknak, és utólag nehezen rekonstruálható a tranzakciók pontos története.

A KasszaSCH célja, hogy erre a problémára egységes, megbízható és könnyen használható digitális elszámolási és nyilvántartási megoldást nyújtson. Ennek megoldásaként a rendszer nyilvántartja a csoportok belső költségvetését, az elérhető termékek készletét, azok ikonjait, támogatja a tagoknak a termékek megvásárlását, tranzakciók hozzárendelését tagokhoz, a tagok számára közösségenként elkülönített belső egyenleget biztosít a költések követésére, és egy adminisztrációs felületet mindezek monitorozására.


\vfill
\selectenglish


%----------------------------------------------------------------------------
% Abstract in English
%----------------------------------------------------------------------------
\chapter*{Abstract}\addcontentsline{toc}{chapter}{Abstract}

Various communities and student groups face a lot of expense-sharing and inventory tasks during their everyday activities, internal events, and camps. They need to track communal supplies, shared expenses, upfront contributions, and group spending. Groups usually handle these accounts on paper, in spreadsheets, or with chat notes. However, these methods are hard to see through, prone to errors, don't let members check their balances in real time, and make it difficult to retrace the exact transaction history later on.

KasszaSCH aims to solve this problem by providing a unified, reliable, and easy-to-use digital accounting and inventory management system. As a solution, the system keeps track of group budgets, available item inventories, and their icons. It also supports member purchases, assigns transactions to specific members, offers separate balances for each community to follow spending, and provides an admin interface to monitor everything.


\vfill
\selectthesislanguage

